
\subsection{Phylogenetic composition}
\label{sec:composition}

The order-1 HMM (Section~\ref{sec:order1hmm}) and transducers (Section~\ref{sec:order1trans})
can be composed on a phylogenetic tree to yield a single composite machine
whose state encodes the joint configuration of all branch machines.

Given a rooted binary tree with $n$ leaves:
\begin{itemize}
\item Number nodes $0,\ldots,2n-2$ in \emph{preorder} (root $= 0$).
\item Place the order-1 Singlet HMM on a notional branch above the root (node 0).
\item Place an order-1 Pair Transducer on each real branch.
\item Each node $v$ carries a \emph{tag} $\in \alphabet \cup \{\varepsilon\}$, initially $\varepsilon$.
\end{itemize}

Each branch machine is either the root HMM or a branch transducer.
All machines are converted to waiting-machine form (they already are, by construction above).

\subsubsection{Composition rules}

\paragraph{State constraint.}
A node's machine may advance (take a transition) only if all higher-numbered nodes' machines
are in waiting states.

\paragraph{Tagging.}
When an internal node $v$'s transducer takes a transition with output symbol $\destok \neq \varepsilon$,
node $v$ is tagged with $\destok$.
If $v$'s tag is non-$\varepsilon$, the next move must feed $v$'s tag as input
to both child branch transducers (forcing them out of their waiting states).
This clears $v$'s tag (but the children's transitions may tag downstream nodes).

\paragraph{Priority.}
If multiple nodes are tagged simultaneously,
the lowest-numbered (closest to root) tagged node is cleared first.

\paragraph{Cascading.}
This system of tags allows upward propagation from observed leaf emissions (MSA columns)
via Felsenstein-style pruning:
working from leaves (known emissions) upward, each internal node's ancestral character
is a latent variable marginalized by the DP.

\subsubsection{Composite state space}

A \emph{configuration} of the composed machine is a tuple
\[
\sigma = (q_0, q_1, \ldots, q_{2n-2},\; \tau_0, \tau_1, \ldots, \tau_{2n-2})
\]
where $q_v$ is node $v$'s machine state and $\tau_v \in \alphabet \cup \{\varepsilon\}$ is its tag.
The start configuration has all machines in $\sta$ and all tags $\varepsilon$.
The end configuration has all machines in $\fin$ and all tags $\varepsilon$.

\paragraph{Practical caveat.}
While this composition defines a valid single machine whose language is the set of MSAs
weighted by the full phylogenetic likelihood, the composite state space is
$O(|Q|^{2n-1} \cdot |\alphabet|^{2n-1})$ where $|Q|$ is the number of states per branch machine---i.e.\
geometric in the number of taxa.
Explicitly constructing the composed machine is therefore impractical for all but
the smallest trees.
The beam algorithms that follow (Sections~\ref{sec:beam-backward}--\ref{sec:progressive})
avoid this by enumerating only the configurations reachable within a pruned beam,
so that the effective state space remains manageable.
The composition formalism is nonetheless useful as a \emph{specification}:
it defines the target distribution from which the beam search samples or
whose expected counts the Forward-Backward algorithm estimates.

\subsection{Beam Backward algorithm (BeamMSA)}
\label{sec:beam-backward}

Given a multiple sequence alignment (MSA) with $L$ columns
and the composite machine from Section~\ref{sec:composition},
we compute the alignment likelihood using a \emph{beam Backward} algorithm,
working from the end configuration backward.

\paragraph{Columns and the emission constraint.}
Each MSA column $\ell = 1,\ldots,L$ specifies, for each leaf $v$,
either a character $y_v^\ell \in \alphabet$ or a gap.
A configuration $\sigma$ is \emph{compatible} with column $\ell$
if the set of leaf emissions implied by $\sigma$'s tags matches the column.

\paragraph{Backward recurrence.}
Let $B(\sigma)$ denote the Backward variable: the total probability of
generating MSA columns $\ell, \ell+1, \ldots, L$ and reaching the end state,
given that the composite machine is currently in configuration $\sigma$ just before column $\ell$.
\begin{align}
B(\sigma_{\fin}) &= 1 && \mbox{(end configuration)} \label{eq:back-base}\\
B(\sigma) &= \sum_{\sigma'} T(\sigma, \sigma')\, B(\sigma') && \mbox{(all other configurations)} \label{eq:back-rec}
\end{align}
where $T(\sigma,\sigma')$ is the composite transition weight
(product of individual machine transitions, subject to the composition rules above)
and the sum is over all successor configurations $\sigma'$.

The alignment likelihood is $B(\sigma_{\sta})$.

\paragraph{Beam pruning.}
Maintain a beam $\mathcal{B}_\ell$ of at most $W$ configurations per column,
ranked by $B(\sigma)$.
When expanding $\mathcal{B}_\ell$ from $\mathcal{B}_{\ell+1}$,
discard any $\sigma$ whose $B(\sigma)$ falls below $B_{\max} / \Delta$
where $B_{\max}$ is the current maximum and $\Delta$ is the beam width ratio.
If the beam collapses (no configurations remain), backtrack.

\paragraph{Epsilon closures within a column.}
A \emph{column-emitting move} is any transition whose output cascades down the tree
to produce a new MSA column: specifically, an insertion (a transition that outputs a character
without consuming input), which then tags the node and cascades to its descendants.
Between column-emitting moves, machines may make silent transitions
(tag propagation, waiting-state transitions, etc.).
These form an $\varepsilon$-closure that must be computed at each step.
For each configuration in the beam, enumerate all reachable configurations
via silent transitions (respecting the priority ordering),
accumulating weights multiplicatively along each path.
In practice, null cycles (if any) can either be ignored (assuming the model's
null-state topology is acyclic) or handled by allowing a configurable number
of extra exploratory steps in the beam search.

\subsubsection{Forward traceback}

After the Backward pass reaches $\sigma_{\sta}$, a stochastic Forward traceback
samples a path from $\sigma_{\sta}$ to $\sigma_{\fin}$:
\begin{align}
P(\sigma' | \sigma) &= \frac{T(\sigma,\sigma')\, B(\sigma')}{B(\sigma)} \label{eq:fwd-sample}
\end{align}
At each step, sample the next configuration proportional to~(\ref{eq:fwd-sample}).
This yields a sampled alignment (including ancestral sequences at internal nodes).

\subsubsection{Beam Forward-Backward}

Alternatively, after the Backward beam pass:
\begin{enumerate}
\item Prune dead-end configurations from each $\mathcal{B}_\ell$
  (those with no predecessor in $\mathcal{B}_{\ell-1}$).
\item Run a Forward pass over the pruned beam:
\begin{align}
F(\sigma_{\sta}) &= 1 \label{eq:fwd-base}\\
F(\sigma') &= \sum_{\sigma \in \mathcal{B}} T(\sigma,\sigma')\, F(\sigma) \label{eq:fwd-rec}
\end{align}
\item Posterior marginals for any feature $\phi$:
\begin{align}
P(\phi | \mbox{MSA}) &= \frac{1}{B(\sigma_{\sta})} \sum_{\sigma \to \sigma'} F(\sigma)\, T(\sigma,\sigma')\, B(\sigma')\, [\phi(\sigma,\sigma')] \label{eq:fb-posterior}
\end{align}
\end{enumerate}

\subsubsection{Beam Viterbi}

Replace $\sum$ with $\max$ in the Backward recurrence~(\ref{eq:back-rec})
and store argmax pointers:
\begin{align}
B^V(\sigma_{\fin}) &= 1 \label{eq:vit-base}\\
B^V(\sigma) &= \max_{\sigma'} T(\sigma,\sigma')\, B^V(\sigma') \label{eq:vit-rec}
\end{align}
The optimal alignment is recovered by Forward traceback following the argmax pointers.

\subsection{Progressive alignment via profile construction (ProgRec)}
\label{sec:progressive}

We now describe a progressive multiple sequence alignment algorithm
using the order-1 machines from Sections~\ref{sec:order1hmm}--\ref{sec:order1trans},
with model parameters from \cite{LargeHolmes2026}.
This follows the transducer-composition approach of
\cite{WestessonEtAl2012},
adapted here for Mealy machines (I/O on transitions rather than states).

The antecedents of this approach are the full multidimensional alignment algorithm of \cite{Hein2000}
which computes the Forward algorithm for TKF91 on a binary tree.
This may be seen as unifying the tree-based Viterbi multiple alignment approach of \cite{Sankoff85}
with the statistical phylogenetics of \cite{Felsenstein81}.
The approach described here also maintains a partial order graph of 
intermediate alignments \cite{LeeGrassoSharlow2002},
which essentially is the approach used by \cite{LoytynojaGoldman2005}.


\subsubsection{Recognizers and Profiles}

Let the phylogenetic tree have $n$ leaves with observed sequences
$\{y_v : v \in \mbox{leaves}\}$, nodes numbered in preorder.
Let $R$ denote the order-1 Singlet HMM (root generator, Section~\ref{sec:order1hmm})
and $B_v$ the order-1 Pair Transducer on the branch to node $v$ (Section~\ref{sec:order1trans}).

\paragraph{Recognizers.}
For each leaf $v$, the \emph{exact-match recognizer} $\mathcal{R}(y_v)$
is a transducer with empty output alphabet that accepts only $y_v$:
it has $|y_v|+1$ states (positions $0,\ldots,|y_v|$),
with a single input-consuming transition $i \to i+1$ labeled by $y_v[i+1]$ at each position.
All states are waiting states except the start.

\paragraph{Profiles.}
A \emph{profile} $E_v$ at node $v$ is a recognizer (empty output alphabet)
that accepts a set of plausible ancestral sequences at $v$,
weighted by their approximate posterior probability given the descendants of $v$.
For leaves, $E_v = \mathcal{R}(y_v)$.

\subsubsection{Progressive reconstruction}

Working from the leaves toward the root, for each internal node $v$ with children $l,r$:

\paragraph{Step 1: Compose branch and profile.}
For each child $c \in \{l,r\}$, form the composition
$B_c \circ E_c$,
which is a transducer mapping the sequence at $v$ to the constrained sequences at child $c$.
Since $E_c$ has empty output, $B_c \circ E_c$ is a recognizer
(it reads a candidate parent sequence and recognizes it with weight proportional to
the probability of generating the descendant data at $c$).

\paragraph{Step 2: Intersect siblings.}
Form the intersection
\[
H_v = (B_l \circ E_l) \cap (B_r \circ E_r)
\]
This recognizer reads a candidate sequence at $v$ and scores it by the joint probability
of both children's descendant data, given that parent sequence.
In the Mealy-machine intersection, the composite state is
$(q_l, e_l, q_r, e_r)$ where $q_c \in B_c$ and $e_c \in E_c$.
Both sides must agree on the same input symbol when both are ready (waiting);
when one side is not waiting, it advances silently while the other stays put.

\paragraph{Step 3: Compose with root prior.}
Form the generator
\[
M_v = R \circ H_v
\]
This has empty input and empty output: it is a weighted automaton
over the empty string, whose total weight $Z = \sum_\pi w(\pi)$ over all paths $\pi$
is the marginal likelihood of the descendant data below $v$
(under the stationary prior $R$ at $v$).

\paragraph{Step 4: Sample paths and construct profile.}
Sample $K$ paths from $P(\pi | M_v) = w(\pi) / Z$ using a Forward pass followed by
stochastic traceback~(\ref{eq:fwd-sample}).
For each sampled path $\pi$, extract the $H_v$-component states visited.
The profile $E_v$ is the sub-recognizer of $H_v$ containing exactly those states
visited by at least $\tau \geq 2$ of the $K$ sampled paths.
This bounds $|E_v| = O(KL)$ where $L = \max_v |y_v|$.

\subsubsection{Mealy-machine composition and intersection}
\label{sec:mealy-ops}

For completeness, we state the composition and intersection rules
for Mealy machines in waiting-machine normal form
(``ready'' states = waiting states with input-consuming transitions only;
``unready'' states = non-waiting, silent transitions only).

\paragraph{Composition.}
Given $T = (\Omega_X, \Omega_Y, \ldots)$ and $U = (\Omega_Y, \Omega_Z, \ldots)$
in Mealy normal form,
$T \circ U$ has states $\subseteq \mathcal{S}_T \times \mathcal{S}_U$ with transition weight:
\[
w''((t,u), \omega_x, \omega_z, (t',u')) = \begin{cases}
\delta_{tt'}\, \delta_{\omega_x\varepsilon}\, w'(u,\varepsilon,\omega_z,u')
& \text{if } u \text{ unready} \\[4pt]
\delta_{uu'}\, \delta_{\omega_z\varepsilon}\, w(t,\omega_x,\varepsilon,t')
+ \displaystyle\sum_{\omega_y} w(t,\omega_x,\omega_y,t')\, w'(u,\omega_y,\omega_z,u')
& \text{if } u \text{ ready}
\end{cases}
\]

\paragraph{Intersection.}
Given $T = (\Omega_X, \Omega_T, \ldots)$ and $U = (\Omega_X, \Omega_U, \ldots)$
in Mealy normal form,
$T \cap U$ has states $\subseteq \mathcal{S}_T \times \mathcal{S}_U$
with output alphabet $\Omega_T \times \Omega_U$ and transition weight:
\[
w''((t,u), \omega_x, (\omega_y,\omega_z), (t',u')) = \begin{cases}
\delta_{tt'}\, \delta_{\omega_x\varepsilon}\, \delta_{\omega_y\varepsilon}\, w'(u,\varepsilon,\omega_z,u')
& \text{if } u \text{ unready} \\[4pt]
\delta_{uu'}\, \delta_{\omega_x\varepsilon}\, \delta_{\omega_z\varepsilon}\, w(t,\varepsilon,\omega_y,t')
& \text{if } t \text{ unready, } u \text{ ready} \\[4pt]
w(t,\omega_x,\omega_y,t')\, w'(u,\omega_x,\omega_z,u')
& \text{if both ready}
\end{cases}
\]

\subsubsection{Forward recursion for $M_v$}
\label{sec:forward-mv}

The generator $M_v = R \circ H_v$ has states $m = (\rho, q_l, e_l, q_r, e_r)$
where $\rho \in R$, $(q_l,e_l) \in B_l \circ E_l$, and $(q_r,e_r) \in B_r \circ E_r$.
The Forward variable $Z(m)$ satisfies:
\begin{align}
Z(\sta_M) &= 1 \label{eq:profile-fwd-base}\\
Z(m') &= \sum_{m : (m,\varepsilon,\varepsilon,m') \in \mathcal{T}} w(m,\varepsilon,\varepsilon,m')\, Z(m) \label{eq:profile-fwd-rec}
\end{align}
where $\mathcal{T}$ is the transition set of $M_v$.
The total likelihood is $Z(\fin_M)$.

The fill order iterates over $e_l$ and $e_r$ in topological order
(corresponding to positions in the child profiles),
with an inner loop over $Q_v = R \circ (B_l \cap B_r)$ states
(the ``comparison kernel'' Pair HMM).
This has time complexity $O(|B|^2\, |E_l|\, |E_r|)$ per internal node,
where $|B|$ is the branch transducer state count.

\subsubsection{Profile extraction}

Given the Forward table $Z$, sample paths $\pi^{(1)},\ldots,\pi^{(K)}$ from $M_v$
using the stochastic traceback~(\ref{eq:fwd-sample}).
For each path $\pi^{(k)}$, let $\mathcal{H}^{(k)} = \{(q_l,e_l,q_r,e_r) : (\rho,q_l,e_l,q_r,e_r) \in \pi^{(k)}\}$
be the $H_v$-states visited.

The profile $E_v$ is the sub-automaton of $H_v$ induced by the states
\[
\mathcal{S}_{E_v} = \{ h \in H_v : |\{k : h \in \mathcal{H}^{(k)}\}| \geq \tau \}
\]
with the same transition weights as $H_v$, restricted to $\mathcal{S}_{E_v}$.
Adding appropriate wait states places $E_v$ in Mealy normal form.

\paragraph{Bubble merging.}
Paths through $E_v$ that traverse the same sequence of wait states
but differ only in latent-state assignments can be merged by collapsing bubbles
(identifying states with identical incoming and outgoing wait-state connectivity).
This further compresses the profile without changing the recognized language.

\subsubsection{MSA extraction}

A sampled path through $M_1$ (the root) determines, at each position,
which $H_1$-state is visited, and therefore which states of $E_l$, $E_r$ are aligned.
Recursing into the child profiles yields a full column assignment for all leaves.

Specifically, each emitting transition in the sampled path implies:
\begin{itemize}
\item a character at the current node (from the root generator or branch match),
\item advancement of the left profile, right profile, or both,
\item and therefore a column in the MSA (with gaps for non-advancing sides).
\end{itemize}

When bubble merging has been applied, the canonical path
(chosen during merging) is used to resolve any ambiguity in the
sub-alignment of the clade below the merged bubble.

\subsubsection{Viterbi Progressive Reconstruction}

An alternative to the sampling-based profile construction of Step~4 above
is to use Viterbi decoding at each internal node.
Instead of sampling $K$ paths from $M_v$ and building a multi-path profile,
the Viterbi variant computes a single maximum-likelihood path through
$M_v$ and uses the resulting ancestral sequence directly as the
reconstructed sequence at node $v$.
This gives a deterministic progressive reconstruction that avoids
the $O(KL)$ profile size but sacrifices the ability to represent
uncertainty in the ancestral sequence.
